\documentclass[a4paper,12pt]{article}
\usepackage[utf8]{inputenc}
\usepackage[T1]{fontenc}
\usepackage[ngerman]{babel}
\usepackage{amsmath,amssymb}
\usepackage{booktabs}
\usepackage{enumitem}
\usepackage[margin=2.5cm]{geometry}
\usepackage{pdflscape}

\begin{document}

\begin{center}
{\LARGE\bfseries \"Ubungsblatt 1 -- L\"osungen}\\[0.5em]
{\large VC Logik}
\end{center}

\bigskip

%% ============================================================
%% AUFGABE 1
%% ============================================================
\subsection*{Aufgabe 1 \textnormal{(1 Punkt)} -- Vollst\"andig geklammerte Formeln}

Operatorpr\"azedenz (von st\"arkster zu schw\"achster Bindung):
$\neg,\; \wedge,\; \vee,\; \to,\; \leftrightarrow$

\begin{enumerate}[(a)]
\item $A \vee \neg\neg B \wedge \neg\neg\neg C
      \;=\; (A \vee (\neg\neg B \wedge \neg\neg\neg C))$

\item $A \vee \neg C \leftrightarrow \neg B \wedge C
      \;=\; ((A \vee \neg C) \leftrightarrow (\neg B \wedge C))$

\item $A \leftrightarrow \neg C \to B \wedge C
      \;=\; (A \leftrightarrow (\neg C \to (B \wedge C)))$

\item $\neg A \leftrightarrow B \to \neg C \wedge A \vee B
      \;=\; (\neg A \leftrightarrow (B \to ((\neg C \wedge A) \vee B)))$
\end{enumerate}

\bigskip

%% ============================================================
%% AUFGABE 2
%% ============================================================
\subsection*{Aufgabe 2 \textnormal{(2 Punkte)} -- Teilformeln}

Formel: $\neg A \leftrightarrow B \to \neg\neg\neg C \wedge A \vee \neg\neg B$

\medskip
Vollst\"andig geklammert:
\[
  (\neg A \leftrightarrow (B \to ((\neg\neg\neg C \wedge A) \vee \neg\neg B)))
\]

Die Formel besitzt \textbf{13~Teilformeln}:

\begin{enumerate}
\item $A$
\item $B$
\item $C$
\item $\neg C$
\item $\neg\neg C$
\item $\neg\neg\neg C$
\item $\neg B$
\item $\neg\neg B$
\item $\neg A$
\item $(\neg\neg\neg C \wedge A)$
\item $((\neg\neg\neg C \wedge A) \vee \neg\neg B)$
\item $(B \to ((\neg\neg\neg C \wedge A) \vee \neg\neg B))$
\item $(\neg A \leftrightarrow (B \to ((\neg\neg\neg C \wedge A) \vee \neg\neg B)))$
\end{enumerate}

\newpage

%% ============================================================
%% AUFGABE 3
%% ============================================================
\begin{landscape}
\subsection*{Aufgabe 3 \textnormal{(2 Punkte)} -- Wahrheitstabelle}

Wahrheitstabelle f\"ur
$\;\neg A \leftrightarrow B \to \neg\neg\neg C \wedge A \vee \neg\neg B$

\medskip
Abk\"urzungen f\"ur die zusammengesetzten Teilformeln:
\begin{align*}
F_{10} &= (\neg\neg\neg C \wedge A) &
F_{11} &= (F_{10} \vee \neg\neg B) &
F_{12} &= (B \to F_{11}) &
F_{13} &= (\neg A \leftrightarrow F_{12})
\end{align*}

\begin{center}
\renewcommand{\arraystretch}{1.25}
\begin{tabular}{ccc|cccccc|cccc}
\toprule
$A$ & $B$ & $C$
  & $\neg C$ & $\neg\neg C$ & $\neg\neg\neg C$
  & $\neg B$ & $\neg\neg B$ & $\neg A$
  & $F_{10}$ & $F_{11}$ & $F_{12}$ & $F_{13}$ \\
\midrule
0 & 0 & 0  & 1 & 0 & 1  & 1 & 0 & 1  & 0 & 0 & 1 & \textbf{1} \\
0 & 0 & 1  & 0 & 1 & 0  & 1 & 0 & 1  & 0 & 0 & 1 & \textbf{1} \\
0 & 1 & 0  & 1 & 0 & 1  & 0 & 1 & 1  & 0 & 1 & 1 & \textbf{1} \\
0 & 1 & 1  & 0 & 1 & 0  & 0 & 1 & 1  & 0 & 1 & 1 & \textbf{1} \\
1 & 0 & 0  & 1 & 0 & 1  & 1 & 0 & 0  & 1 & 1 & 1 & \textbf{0} \\
1 & 0 & 1  & 0 & 1 & 0  & 1 & 0 & 0  & 0 & 0 & 1 & \textbf{0} \\
1 & 1 & 0  & 1 & 0 & 1  & 0 & 1 & 0  & 1 & 1 & 1 & \textbf{0} \\
1 & 1 & 1  & 0 & 1 & 0  & 0 & 1 & 0  & 0 & 1 & 1 & \textbf{0} \\
\bottomrule
\end{tabular}
\end{center}

\medskip
Die Formel ist genau dann wahr, wenn $A = 0$ gilt (d.\,h.\ sie ist \"aquivalent zu $\neg A$).

\end{landscape}

\end{document}
